\section{Outer circuits without parity gates}
For every fixed $\eta>0$ there is a fixed $c<1$ such that recognizing
$A_{\lfloor cN/n\rfloor}$ with pointwise error at most $1/3$ requires
at least $(13/4-\eta)N$ outer $U_2$ gates for all sufficiently large $n$.
Here $U_2$ consists of the binary Boolean gates except XOR and its
complement, with free wire negations. The inner circuit convention and
the easy input law are those of Theorem~\ref{th:law}. Consequently the
threshold-supplied problem requires $(13/4-o(1))N$ gates in this outer basis.

To see the relevant simulation explicitly, use the projection recursion
of Golovnev, Kulikov, Smal and Tamaki (\href{https://golovnev.org/papers/sharp_sat.pdf}
{full version, Lemma 7 and Corollary 8(1)}). Write $s$ for the residual
number of gates, $i$ for the number of live input variables, and $i_1$
for the number used once. For the nonnegative measure
$M=s+a i-i_1/2$, their decreases at the children of each nonterminal
split dominate one of the four vectors
\[
 (a,2a),\quad(2a,2a,2a,3a),\quad
 (a+7/2,a+3),\quad (a+9/2,a+2).
\]
A split fixes one or two variables, either to constants or to signed
copies of remaining variables. Each split is a parity decision subtree
of depth at most two. The substitution and normalization compute the
actual restricted function on each branch. When the residual output is
trivial, at most a constant number of extra parity queries determines
it. Thus bounding the number of leaves of this recursion bounds an
exact representation of the original Boolean function.

For completeness, set $\lambda=2^{1/(a+13/4)+K/a^3}$ for an absolute
constant $K$ and sufficiently large $a$. For each vector $v$ above,
$\sum_{u\in v}\lambda^{-u}\le1$. For the first two vectors the sums
tend to $3/4$ and $7/8$, respectively. Each of the last two vectors
has entries $a+13/4\pm d$ for bounded $d$, so its sum is
$2\lambda^{-(a+13/4)}\cosh(d\ln\lambda)$, at most one upon taking $K$
sufficiently large. Induction on the tree bounds its leaf count by
$O(\lambda^{M})$. If initially $s\le(13/4-\eta)N$, the base-two
logarithm of this bound is at most
\[
 N\left(1-\frac{\eta}{a+13/4}+O(a^{-2})\right)+O(1).
\]
Fixing $a$ sufficiently large in terms of $\eta$ gives at most
$2^{(1-\kappa_\eta)N}$ leaves for a constant $\kappa_\eta>0$.
Choose $1-\kappa_\eta<\alpha<c<1$ and apply the parity leaf lower
bound to each fixing of the outer circuit's random bits.
